% To analysize the information composition of different speech tokens quantitatively, we established Discrete Speech Representation Information Benchmark (DSRIBench) based on SUPERB~\citep{yang2021superb}. In this work, we mainly investigate two speech attributes: content and timbre. 

 % We calculate the mutual information between speech tokens and different attributes through the performance of downstream tasks .

To build powerful speech language models, discrete speech representations should possess the following two key characteristics: i)~Strong alignment with text; ii) Effective preservation of speech information. Building on this premise, we establish speech Language Model Token Benchmark~(\bench) to assess the suitability of speech tokens for constructing speech language models. 
% In this benchmark, we evaluate the alignment between speech tokens and text by estimating their mutual information. We assess preservation of speech information within speech tokens by evaluating the quality of resynthesized speech. 
% \bench is available on \url{https://github.com/0nutation/SLMTokBench}.

\subsection{Text Alignment Evaluation}
\label{sec:026_information principle}

% Mutual Information is a measure of how much the knowledge of one variable reduces the uncertainty about another variable. It quantifies the level of dependence or relationship between the variables.
We evaluate the degree of text alignment by estimating the mutual information between speech tokens and text. 
For notation, $\mathbf{X}$ denotes discrete speech representations; $\mathbf{Y}$ denotes text; $\mathcal{I}(\mathbf{X};\mathbf{Y})$ denotes the mutual information; test dataset is denoted as $\mathcal{D} =\{(x_i, y_i)\}_{i=1}^{N}$ and $\theta$ denotes the downstream model. Through the derivation in Appendix~\ref{sec:app:mi_estimation}, we can estimate
$\mathcal{I}(\mathbf{X};\mathbf{Y})$ as:
\begin{align*}
    \mathcal{\hat{I}}(\mathbf{X};\mathbf{Y})=\frac{1}{N^2}\sum_{i=1}^{N}\sum_{j=1}^{N}[\log q_{\theta}(y_i|x_i)-\log q_{\theta}(y_j|x_i)]
\end{align*}
where $q_{\theta}(\mathbf{Y}|\mathbf{X})$ is the variational distribution and can be parameterized by the downstream model $\theta$.

% A measure of mutual information between variable $\mathbf{X}$ and $\mathbf{Y}$ can be formulated as:
% \begin{align}
%     \mathcal{I}(\mathbf{X};\mathbf{Y})=\int_{\mathbf{X}}\int_{\mathbf{Y}}log\frac{P(\mathbf{X},\mathbf{Y})}{P(\mathbf{X})P(\mathbf{Y})} 
% \end{align}
% where $P(\mathbf{X})$ and $P(\mathbf{Y})$ are the marginal distributions of $\mathbf{X}$ and $\mathbf{Y}$
%  respectively, and $P(\mathbf{X},\mathbf{Y})$ denotes the joint distribution of \mathbf{X} and \mathbf{Y}.
 
% The variational contrastive log-ratio upper bound~(vCLUB)~\citep{cheng2020club} of mutual information is defined by:
% \begin{align}
%     \mathcal{I}(\mathbf{X};\mathbf{Y})=E_{p(\mathbf{X},\mathbf{Y})}[\log q_{\theta}(\mathbf{Y}|\mathbf{X})]
%     -E_{p(\mathbf{X})p(\mathbf{Y})}[\log q_{\theta}(\mathbf{Y}|\mathbf{X})]
% \end{align}
% where $q_{\theta}(\mathbf{Y}|\mathbf{X})$ is the variational distribution to approximate the ground-truth probability $P(\mathbf{Y}|\mathbf{X})$ and can be parameterized by the downstream model $\theta$.

% With test dataset $\mathcal{D}$, $\mathcal{I}(\mathbf{X};\mathbf{Y})$ has an unbiased estimation as:
% \begin{align}
%     \mathcal{\hat{I}}(\mathbf{X};\mathbf{Y})=\frac{1}{N^2}\sum_{i=1}^{N}\sum_{j=1}^{N}[\log q_{\theta}(y_i|x_i)-\log q_{\theta}(y_j|x_i)]
% \end{align}
The downstream model is a vanilla 2-layer 1024-unit BLSTM optimized by CTC loss on characters and it takes speech tokens as inputs. Specifically, for each discrete representation, we first establish an embedding matrix, which can be either randomly initialized or derived from the k-means centroid matrix or vector quantization codebooks obtained during the discretization process. We use the embedding matrix to embed the discrete representations and obtain continuous representations, which are then fed into the downstream models. We train the downstream model on LibriSpeech train-clean-100 subset and use dev-clean subset for estimating mutual information. We also calculate the word error rate~(WER) on the test set.
For downstream model training, we configure the training setup with a batch size of 32, a learning rate of 1e-4, and a total of 200k global steps.


% This section explains how to compute mutual information using the performance of downstream tasks.
% For notation, $\mathbf{X}$ denotes discrete speech representations; $\mathbf{Y}$ denotes one of the three speech attributes: content, timbre, or paralinguistics; the test dataset is denoted as $\mathcal{D} =\{(\mathbf{x}, \mathbf{y})\}; $$H(\mathbf{X})$ denotes the entropy of $\mathbf{X}$; $H(\mathbf{Y})$ denotes the entropy of $\mathbf{Y}$; $H(\mathbf{Y}|\mathbf{X})$ denotes the entropy of $\mathbf{Y}$ conditional on $\mathbf{X}$; $H(\mathbf{X}|\mathbf{Y})$ denotes the entropy of $\mathbf{X}$ conditional on $\mathbf{Y}$; $I(\mathbf{Y};\mathbf{X})$ denotes the mutual information; $N$ denotes the number of samples in the test dataset.

% Because each sample is randomly sampled from test dataset, we get
% \begin{align}
%     p(x)&=\frac{1}{N}, \forall x \in D
% \end{align}
% We have
% % \begin{equation}
% \begin{align}
% I(\mathbf{X};\mathbf{Y})&=H(\mathbf{Y})-H(\mathbf{Y}|\mathbf{X}) \\
% H(\mathbf{Y})&=-\sum_{y \in D}p(y)\log p(y) \\
% H(\mathbf{Y}|\mathbf{X})&=-\sum_{(x,y) \in D}p(x,y)\log p(y|x)\\
% &=-\sum_{(x,y) \in D}p(y|x)p(x)\log p(y|x)\\
% &=-\frac{1}{N}\sum_{(x,y) \in D}p(y|x)\log p(y|x)
% \end{align}
% % \end{equation}


% So
% \begin{align}
% I(\mathbf{X};\mathbf{Y})=\frac{1}{N}\sum_{(x,y) \in D}p(y|x)\log p(y|x)-\sum_{y \in D}p(y)\log p(y)
% \end{align}
% where for content, $p(y)$ can be calculated from large language model like GPT-3; for timbre and paralinguistics, $p(y)$ can be calculated from dataset statistics. $p(x|y)$ can be calculated from downstream model output.


\subsection{Information Preservation Evaluation}
\label{sec:027_bench_resys}
To evaluate the preservation of speech information in discrete speech representations, we convert speech tokens back to speech and evaluate resynthesized speech by automatic metrics on content and timbre.
We train a unit-HiFIGAN~\citep{polyak2021speech} on LibriSpeech dataset to convert HuBERT units to waveform. Notably, to avoid interference from additional information, we don't supply any speaker information during training. For Encodec tokens, we used the Encodec decoder to
directly produce the waveform.
Content preservation is evaluated by computing the WER through transcribing the resynthesized speech using the Whisper en-medium model~\citep{radford2023robust}. Timbre preservation is evaluated by utilizing WavLM-TDNN~\citep{Chen_2022} to calculate speaker similarity between the synthesized and groundtruth speech. We randomly sample 300 speech samples from LibriSpeech test set for evaluation. 


\subsection{Comparing Semantic $\&$ Acoustic Tokens}
\label{sec:028_bench_results}
We use HuBERT L9 units to represent semantic tokens and EnCodec codes to represent acoustic tokens.
As shown in Table~\ref{tab:bench_results}, semantic tokens achieve high mutual information with text but their resynthesized speech has low speaker similarity. 
% That means semantic tokens align well with the text and are excellent for capturing the meaning and context of spoken words, but they can lack the paralinguistic features of speech, such as timbre, or prosody.
Acoustic tokens achieve low WER and high speaker similarity for resynthesized speech but have low mutual information with text.
% That shows acoustic tokens do tend to capture the rich details of the speech signal, such as pitch, timbre, but they often lack the semantic interpretability that would align them closely with the corresponding text. 


% We use HuBERT L9 units to represent semantic tokens and EnCodec codes to represent acoustic tokens.
% As shown in Table~\ref{tab:bench_results}, semantic tokens achieve high mutual information with text but their resynthesized speech has low speaker similarity. 
% That means semantic tokens align well with the text and are excellent for capturing the meaning and context of spoken words, but they can lack the paralinguistic features of speech, such as timbre, or prosody.
% Acoustic tokens achieve low WER and high speaker similarity for resynthesized speech but have low mutual information with text.
% That shows acoustic tokens do tend to capture the rich details of the speech signal, such as pitch, timbre, but they often lack the semantic interpretability that would align them closely with the corresponding text. 

% % \begin{table*}[h]
%     % \setlength{\abovecaptionskip}{-0.cm}
%     \setlength{\belowcaptionskip}{-0.2cm}
%     % \renewcommand{\arraystretch}{1.2}
%     \Large
%     \centering
%     \resizebox{0.7\textwidth}{!}{
%     \begin{tabular}{l|cc|cc|cc}
%         \toprule
%         ~ & \multicolumn{2}{c}{\textbf{Content}} & \multicolumn{2}{c}{\textbf{Timbre}}  & \multicolumn{2}{c}{\textbf{Paralinguistics}} \\
%         ~ & MI & WER & MI& ACC & MI & ACC \\
%         \midrule
%         Semantic tokens &31.2 &9.88 & 9.9& 0.66 & 3.3 & 0.50 \\
%         Acoustic tokens~(RVQ-1) &16.5 &61.52 &6.52 &0.44 &1.8 & 0.50\\
%         Acoustic tokens~(RVQ-8) & 23.6 & 30.91 & 14.7 & 0.94 & 1.6 & 0.49 \\
%         \bottomrule
%     \end{tabular}}
%     \caption{Results of semantic tokens and acoustic tokens on DSRIBench.
%     }
%     \label{tab:disrbench_results}
% \end{table*}

\begin{table*}[h]
    % \setlength{\abovecaptionskip}{-0.cm}
    \setlength{\belowcaptionskip}{-0.2cm}
    % \renewcommand{\arraystretch}{1.2}
    \Large
    \centering
    \resizebox{0.6\textwidth}{!}{
    \begin{tabular}{lc|cc|cc}
        \toprule
        ~ & & \multicolumn{2}{c}{\textbf{TA}} & \multicolumn{2}{c}{\textbf{IP}}   \\
        ~ & & MI & WER$^{\dagger}$ & WER$^{\ast}$ & SIM\\
        \midrule
        Semantic tokens & & \textbf{31.2} & \textbf{9.88} & & \\
        Acoustic tokens & RVQ-1 & 16.5 & 61.52 & 38.34 & 0.92 \\
        Acoustic tokens & RVQ-8 & 23.6 & 30.91 & \textbf{5.11} & \textbf{0.98} \\
        \bottomrule
    \end{tabular}}
    \caption{Results of semantic tokens and acoustic tokens on DSRIBench. TA refers to text alignment and IP denotes information preservation. MI and WER$^{\dagger}$ refer to mutual information and word error rate of the downstream model. WER$^{\ast}$ and SIM refer to word error rate and speaker similarity of resynthesized speech respectively.
    }
    \label{tab:disrbench_results}
\end{table*}




